\documentclass{ip-quiz}

\quizlevel{n2}
\quiznum{1}
\quiztitle{Funciones y strings}

\begin{document}

\makeheader
\maketitle

\begin{sectionbox}{Enunciado}
Un sensor de velocidad en Bogotá captura placas vehiculares con baja precisión.

Dada una placa fotografiada y tres placas candidatas, determine cuál candidata es más similar a la capturada por el sensor. La similitud vale \texttt{1} si las letras coinciden y \texttt{1} si los dígitos coinciden (entre \texttt{0} y \texttt{2} en total). En caso de empate, retorne la última candidata.

Retorne un mensaje que indique claramente cuál placa ganó.

\textbf{Nota: } Las placas en Colombia constan de tres letras, un espacio y tres números. E.g., \texttt{ABC 123}.
\end{sectionbox}

\pagebreak

\begin{sectionbox}{Solución}
\begin{minted}[escapeinside=||]{python}
def calcular_similitud(placa_foto: str, placa_x: str) -> int:
    puntaje = 0
    if placa_foto[:3] == placa_x[:3]:
        puntaje += 1
    |\tms{ifD}|if|\tme{ifD}| placa_foto[4:] == placa_x[4:]:
        puntaje += 1
    return puntaje


def encontrar_placa_mas_similar(
    |\tikzmark{ptop}|placa_foto: str,|\tikzmark{pright}|
    placa_1: str,
    placa_2: str,
    placa_3: str|\tikzmark{pbot}|
) -> str:
    mejor_nombre = "placa 1"
    mejor_placa = placa_1
    mayor_similitud = calcular_similitud(placa_foto, placa_1)

    similitud_2 = calcular_similitud(placa_foto, placa_2)
    |\tms{ifA}|if|\tme{ifA}| similitud_2 |\tms{geqA}|>=|\tme{geqA}| mayor_similitud:
        mejor_nombre = "placa 2"
        mejor_placa = placa_2
        mayor_similitud = similitud_2

    similitud_3 = calcular_similitud(placa_foto, placa_3)
    |\tms{ifB}|if|\tme{ifB}| similitud_3 |\tms{geqB}|>=|\tme{geqB}| mayor_similitud:
        mejor_nombre = "placa 3"
        mejor_placa = placa_3
        mayor_similitud = similitud_3

    |\tikzmark{rettop}|return (f"La placa más similar es la {mejor_nombre}.\n" +
            f"Su valor es: {mejor_placa}.\n" +
            f"Su puntaje de similitud es: {mayor_similitud}.")|\tikzmark{retbot}|
\end{minted}

\begin{tikzpicture}[remember picture, overlay]

  \node[box, fit=(ifDs)(ifDe)] (cifD) {};
  \node[note, anchor=north west, text width=5cm]
    (noteA) at ([xshift=-8.9cm] current page.east |- ifDs)
    {Uso \texttt{if} en lugar de \texttt{elif} porque de lo contrario si las letras coinciden, no se evalúan los dígitos.};
  \draw[arr, dotted] (cifD.east) to[bend right=5] (noteA.west);

  \codebrace{pright}{15pt}{ptop}{pbot}{3pt}{5cm}
    {Pongo los parámetros en varias líneas para no rebasar el margen. Igual los separo por comas, como siempre.}

  \node[box, fit=(geqAs)(geqAe)] (cgeqA) {};
  \node[box, fit=(geqBs)(geqBe)] (cgeqB) {};
  \node[note, anchor=west, text width=3.8cm]
    (noteC) at ([xshift=-6.9cm] current page.east |- geqAs)
    {>= y no >: en empate se actualiza al último candidato, como pide el enunciado.};
  \draw[arr, dotted] (cgeqA.east) to[bend left=5]  (noteC.west);
  \draw[arr, dotted] (cgeqB.east) to[bend right=10] (noteC.west);

  \node[box, fit=(ifBs)(ifBe)] (cifB) {};
  \node[note, anchor=north west, text width=5.7cm]
    (noteB) at ([xshift=-8.9cm, yshift=-4pt] current page.east |- ifBs)
    {Similar a arriba, si se usara \texttt{elif} en lugar de \texttt{if}, el código fallaría cuando la placa 3 es la mejor y la 1 la peor: se retornaría la placa 2.};
  \draw[arr, dotted] (cifB.east) to[bend right=5] (noteB.west);

  \codebrace{retbot}{12pt}{rettop}{retbot}{2pt}{3.2cm}
    {Concatenación con + y \textbackslash n dentro de paréntesis para poder escribir el \texttt{str} en varias líneas.}

\end{tikzpicture}
\end{sectionbox}

\end{document}
